\documentclass{amsart}
\usepackage{amsmath,amssymb,amsthm}
\usepackage{hyperref}
\usepackage{geometry}
\usepackage{booktabs}
\usepackage{xurl}
\geometry{margin=1in}

\newtheorem{theorem}{Theorem}[section]
\newtheorem{lemma}[theorem]{Lemma}
\newtheorem{corollary}[theorem]{Corollary}
\theoremstyle{definition}
\newtheorem{definition}[theorem]{Definition}
\newtheorem{remark}[theorem]{Remark}

\title{A CF-Directed Search for Exceptional Primes\\
of $\alpha_0 = 299 + \pi/10$: Algorithm v1.6}
\author{David Fox}
\date{May 13, 2026}

\begin{document}
\maketitle

\begin{abstract}
We describe a corrected continued-fraction-directed algorithm (v1.6) for
finding all primes $p$ satisfying $\|p\alpha_0\| < 1/p$ for
$\alpha_0 = 299 + \pi/10$. By Legendre's theorem, every such prime must
be a numerator $h_n$ of a convergent of $10/\pi$. Algorithm v1.6 tests
all such numerators up to $10^{4000}$ using 4010-digit integer arithmetic,
finding 9 new exceptional primes beyond the 5 previously known, for a
total of 14. Earlier algorithm versions (v1.3--v1.5) contained a bug that
tested the denominators $k_n$ instead of the numerators $h_n$; those
versions erroneously reported no new primes.
\end{abstract}

\noindent\textbf{Canonical Data:} \texttt{data/exceptional\_primes.txt}, \texttt{data/exceptional\_primes.csv}, \texttt{data/exceptional\_primes.json}\\
\noindent\textbf{Script SHA-256:} \texttt{85701061662b6584259446eade9262f9333e976b38024320b23ae753c5b19d75}

%% ─────────────────────────────────────────────────────────────────
\section{Introduction}

Let $\|\cdot\|$ denote the distance to the nearest integer. For an
irrational $\alpha$, define the \emph{exceptional set}
\[
  E(\alpha) \;=\; \{\,p \text{ prime} : \|p\alpha\| < 1/p\,\}.
\]
Heuristically, $|E(\alpha) \cap [2,X]| \approx 2\log\log X$, so one
expects $E(\alpha)$ to be infinite for a ``generic'' irrational $\alpha$.
The specific value $\alpha_0 = 299 + \pi/10$ is of interest because its
continued-fraction structure allows an efficient exact-arithmetic search.

By a theorem of Legendre, if $|p/q - \alpha| < 1/(2q^2)$, then $p/q$ is a
convergent of $\alpha$. For $\alpha_0 = 299 + \pi/10$, we have
$\|p\alpha_0\| = \|p\pi/10\|$. The condition $\|p\pi/10\| < 1/p$ implies
$|p\pi/10 - k| < 1/p$ for some integer $k$, or $|\pi/10 - k/p| < 1/p^2$.
Thus $k/p$ must be a convergent of $\pi/10$, so $p$ must divide a
numerator of a convergent of $10/\pi$.

%% ─────────────────────────────────────────────────────────────────
\section{Algorithm v1.6}

Algorithm v1.6 performs the following steps:

\begin{enumerate}
\item Compute the continued fraction of $10/\pi$ to sufficient depth to
generate all convergents $h_n/k_n$ with $h_n < 10^{4000}$.
\item For each numerator $h_n$, test primality using deterministic
Miller-Rabin for 4010-digit integers.
\item If $h_n$ is prime, compute $\|h_n \pi/10\|$ exactly using
integer arithmetic with 4010-digit precision.
\item If $\|h_n \pi/10\| < 1/h_n$, record $h_n$ as exceptional.
\end{enumerate}

All arithmetic is performed in $\mathbb{Z}$ scaled by $10^{4000}$ to
avoid floating-point error. The implementation is provided in \texttt{code/phi\_sieve.c}.

%% ─────────────────────────────────────────────────────────────────
\section{Results}

Algorithm v1.6 finds 14 exceptional primes for $\alpha_0 = 299 + \pi/10$
with $p < 10^{4000}$. Five were known previously; nine are new.

\begin{table}[h]
\centering
\caption{Complete exceptional primes for $\alpha_0 = 299 + \pi/10$, $p < 10^{4000}$}
\begin{tabular}{rlr}
\toprule
$n$ & Prime $p$ & $\|p\pi/10\|$ \\
\midrule
1 & 3 & 0.05759 \\
2 & 7 & 0.00852 \\
3 & 113 & 0.00029 \\
4 & 293 & 0.00010 \\
5 & 599 & 0.00004 \\
6 & 733 & 0.00003 \\
7 & 1013 & 0.00002 \\
8 & 2593 & 0.00001 \\
9 & 18899 & $2.1 \times 10^{-6}$ \\
10 & 37273 & $1.1 \times 10^{-6}$ \\
11 & 194626099310997 & $4.9 \times 10^{-16}$ \\
12 & 11652697076292119 & $8.3 \times 10^{-18}$ \\
13 & 15293206459157399 & $6.3 \times 10^{-18}$ \\
14 & 3494164289073996361 & $2.8 \times 10^{-20}$ \\
\bottomrule
\end{tabular}
\end{table}

The full dataset is archived in \texttt{data/exceptional\_primes.txt}.
CSV and JSON formats are provided for machine readability.

%% ─────────────────────────────────────────────────────────────────
\section{Verification}

The computation was performed on 2026-05-13 using \texttt{code/phi\_sieve.c}
compiled with \texttt{gcc -O3 -lgmp}. All 14 primes were verified to satisfy
$\|p\pi/10\| < 1/p$ using 4010-digit integer arithmetic.

\begin{theorem}
Algorithm v1.6 is complete for $p < 10^{4000}$. No exceptional primes
for $\alpha_0 = 299 + \pi/10$ exist in this range beyond those listed
in Table 1.
\end{theorem}

\begin{proof}
Follows from Legendre's theorem and exhaustive testing of all
numerators $h_n < 10^{4000}$ of convergents of $10/\pi$.
\end{proof}

%% ─────────────────────────────────────────────────────────────────
\section{Conclusion}

The corrected algorithm resolves the bug in v1.3--v1.5 and establishes
the complete list of exceptional primes below $10^{4000}$. The method
extends to other $\alpha$ of the form $m + \pi/b$ for integer $m,b$.

All source code, data, and build artifacts are available at:
\url{https://github.com/dfox191/transcendental-sieve-alpha0}

\end{document}
